% ============================================================
%  When the Machine Stops Looking
%  Autonomous Replenishment, Censored Demand, and the
%  Vanishing Value of a Human Filter
%
%  TERMINOLOGY LOCK (one term per object, used everywhere):
%    - "autonomous policy"      : the unattended machine learner
%    - "censoring-aware"        : de-censors via Kaplan-Meier
%    - "contamination-blind"    : cannot flag corrupted records
%    - "human filter"           : domain expert who screens data
%    - "observable boundary"    : largest order recently placed,
%                                 above which demand is unidentified
%    - "biased equilibrium"     : attracting order level q_infty < q*
%    - "exposure law"           : d[service]/d[eps] = -(1-beta)
%    - "correction-log alarm"   : minimal filter, fires on logged correction days
%
%  LOCAL STYLE (enforced throughout):
%    - no em dashes, no semicolons, no emphasis font
%    - unpack definitions with "that is" / "namely"
%    - lead with the managerial question, keep the benchmark visible
%    - after each result: math statement, mechanism, managerial meaning,
%      failure mode
% ============================================================

\documentclass[11pt]{article}

\usepackage[margin=1in]{geometry}
\usepackage{amsmath,amssymb,amsthm}
\usepackage{mathtools}
\usepackage{microtype}
\usepackage{enumitem}
\usepackage{xcolor}
\usepackage{mdframed}
\usepackage{booktabs}
\usepackage{titlesec}
\usepackage[round]{natbib}
\bibliographystyle{plainnat}
\usepackage[hidelinks]{hyperref}

\theoremstyle{plain}
\newtheorem{theorem}{Theorem}
\newtheorem{lemma}{Lemma}

\titlespacing*{\section}{0pt}{1.4em}{0.6em}
\setlength{\parskip}{0.35em}
\linespread{1.08}

\title{\vspace{-2.5em}\textbf{When the Machine Stops Looking:\\
Autonomous Replenishment, Censored Demand,\\
and the Vanishing Value of a Human Filter}}
\author{}
\date{}

\begin{document}
\maketitle
\vspace{-3em}

\begin{abstract}
\noindent
Retailers increasingly delegate replenishment to algorithms that learn from
sales with no human in between. A learner can correct for demand censoring on its
own when stockouts are recorded on clean days, but it cannot flag contaminated
records without domain knowledge or access to the returns-and-corrections log.
We study a censoring-aware but contamination-blind learner in a repeated
newsvendor in which a fraction of high-demand events is masked by returns or
corrections booked to those periods, making them appear as low-demand days.
The order and the data then form a closed
loop. We show that the loop drives the order to a unique equilibrium below the
optimum, that no tipping point arises along the way, and that the marginal service
loss per unit of contamination equals the item's stockout frequency $1-\beta$. The
items a firm automates first, namely low-margin, high-volume, thin-service items,
degrade fastest. The loss accrues every period, so the value of a human who
screens the data grows with the planning horizon and is concentrated in screening
the data rather than in overriding orders, which matches the field evidence on
human discretion. A minimal filter that audits only the returns-and-corrections log
recovers nearly first-best service. We validate the equilibrium, the exposure law,
and the filter on a calibrated instance.
\end{abstract}

\section{Introduction}

% P1 Context. Realistic regime first. Names every decision object.
Large retailers now let software place replenishment orders with no buyer in
between. The system reads recent sales, updates its view of demand, and sets the
order quantity $q_t$ for each of hundreds of thousands of items, night after
night. A buyer used to do two things that the order quantity alone does not
record. She read a sold-out day as understating true demand, and she set aside the
records she judged unreliable, that is, the week a promotion, a returns batch, or
a feed error distorted the numbers. The first correction is mechanical when
labels are clean. The firm knows when an item stocked out, so a machine can make
it. The second correction needs to know the business, and it matters most when
returns, phantom inventory, or feed errors mask a stockout in the feed. This paper
asks what the firm loses when the second correction goes away, and shows that for
a recognizable class of items the loss is not a steady efficiency tax but a
service level that drifts down and settles well below target.

% P2 Trade-off. One concrete numeric vignette, both extremes, neither works.
Consider an item the firm wants stocked to a fifty percent service level, so that
the cost-minimizing order is the median of demand, namely $100$ units. Suppose
that on a tenth of the high-demand events, a returns batch or data correction is
booked to the same period, reducing the recorded demand so that what should be
recorded as a stockout instead appears as an ordinary low-sales day, say around
$40$ units on average. A censoring-aware learner corrects for the stockouts it can
see, yet it reads these masked days at face value as evidence that demand is weak,
so it trims the order. The lower order produces more high-demand events, more of
which can be masked by coincidental returns or corrections, and the order keeps
drifting down until it settles near $96$ units at a realized service level of
forty-four percent. Raise the masking rate to thirty percent of high-demand days
and the same item settles near $83$ units at twenty-nine percent service, and a
low-service item that started at thirty percent settles in single digits. The two
clean alternatives both fail. Full automation is cheap per decision and reacts
overnight, but it walks the item down. Full human review corrects every record by
checking the returns log and removing the booked offsets, but it spends scarce
attention on every item the firm carries, and the field evidence is that letting
buyers also override the order itself reduces profit \citep{kesavan2020override}.
The question is where between these extremes the firm should sit, and the answer
turns on how much of the high-demand events are masked by returns and on how
exposed the item is to the loop.

% P3 Difficulty. Why hard on its own terms. No answer preview.
The difficulty is that the order and the data form a closed loop. Because
stockouts are recorded on clean days, one might expect a censoring-aware learner
to be safe, and on clean records it is. The catch is that the de-censoring
correction trusts the stockout labels, and contamination masks exactly those
labels. A masked stockout enters the learner's data as a genuine low-demand day,
lowering the order, and the lower order raises the share of days that stock out
and so can be masked. The learner cannot separate a genuinely weak demand signal
from a corrupted one, because the evidence that would separate them, namely the
demand it failed to serve, never appears in the feed. Saying where this loop
settles, and how fast it descends, means tracking a decision-dependent estimator
together with the order it feeds, and neither piece can be analyzed in isolation.

% P4 Research Question. Closest prior work + restrictive assumption + interrogative
% question + method shift. No answer preview.
The closest prior work studies learning under censoring on clean data.
\citet{besbes2013implications} isolate the regret cost of censoring,
\citet{lariviere1999stalking} show that ordering above the myopic level keeps the
upper tail visible and escapes the censoring trap, and \citet{ding2024feature}
build censoring-aware learners that attain vanishing regret with features.
\citet{hssaine2024datadriven} make the observable boundary central, but in an
offline setting with uncorrupted records. \citet{mersereau2015censored} show that
inventory record inaccuracy compounds the downward bias already present under
censoring. Each of these assumes the records are clean, or treats inaccuracy
separately from the learning loop, which is the one correction a machine cannot
make on its own, because contamination, unlike censoring, is not self-revealing
in the demand feed. So the question this paper asks is the following. When a
censoring-aware but contamination-blind learner runs unattended, does data
contamination wash out as noise, or does it compound through the censoring loop
into a self-reinforcing decline, and how little human screening of the data is
enough to stop it? Answering it requires coupling a contamination model with an
endogenous observable boundary inside a repeated newsvendor, and then pricing a
scarce human filter against the loop.

% P5 Managerial answer. Plain language, no theorem sentences, points forward.
Contamination does not wash out. It compounds. The order descends to a single
equilibrium below the target and stays there, and the descent is smooth, so there
is no warning tipping point to watch for. What governs the damage is one number.
The service the firm loses per unit of contamination equals the item's stockout
frequency, namely one minus its target service level. Low-service, thin-margin,
high-volume items run with high stockout frequencies, so they bleed service
fastest, and they are the items a firm automates first. Because the shortfall
repeats every period, the cost of leaving the loop unattended grows with the
planning horizon rather than washing out, and the value of a human sits entirely
in screening the data rather than in overriding the order, which is why
input-level human corrections help in the field while order overrides do not
\citep{kesavan2025judgmental}. The remedy is cheap and pointed. Masked stockouts
leave a trace in the returns-and-corrections log, so a human called in only on
logged correction days audits exactly the records at fault, and because the cost
penalty grows with the square of the contamination, auditing even half of those
days removes most of it while auditing all of them restores the target. The
contributions below make each of these statements precise and validate them on a
calibrated instance.

% Contributions block. Formal verbs, explicit orders, result references.
\paragraph{Contributions.}
\begin{enumerate}[leftmargin=1.5em,itemsep=0.35em,label=\textbf{\arabic*.}]
  \item \textbf{Model.} We formulate the repeated newsvendor as a closed loop in
  which a censoring-aware but contamination-blind learner sets orders, demand is
  censored at the order, and a fraction $\epsilon$ of true stockouts is masked in
  the feed as low-sales days drawn from a downward source $G$, following the
  returns-and-inventory-inaccuracy patterns documented by
  \citet{mersereau2015censored,dehoratius2008retail}. Lemma~\ref{lem:mf} derives
  the population estimate the learner converges to, namely
  $\widehat F_q(x)=F(x)+\epsilon(1-F(q))G(x)$, in which the contamination term
  carries the stockout frequency $1-F(q)$ and so depends on the order.

  \item \textbf{Main result, with a negative result inside it.}
  Theorem~\ref{thm:main} is the spine of the paper. Part~(i) proves that the
  projected Robbins--Monro learner converges almost surely to a unique equilibrium
  $q_\infty(\epsilon,\beta)<q^\star$ and that the order map is monotone, so no
  tipping point or bistable collapse can arise from this contamination. Part~(ii)
  is the centerpiece, the exposure law
  $\mathrm{d}\,[\text{service}]/\mathrm{d}\epsilon=-(1-\beta)G(q^\star)$, equal to
  $-(1-\beta)$ in the relevant range, so degradation is steepest for low-service
  items. Part~(iii) shows long-run average cost exceeds the optimum by
  $\Theta\!\big((b+h)\,\epsilon^2(1-\beta)^2/f(q^\star)\big)$.

  \item \textbf{Prescription.} Because that penalty is quadratic in contamination,
  a correction-log alarm that audits a fraction $\varphi$ of logged correction days
  removes more than three quarters of it at half coverage and restores the optimum
  at full coverage (Theorem~\ref{thm:main}(iii), Table~\ref{tab:sim}). Oversight
  should be triggered by the returns and corrections log rather than by stockout
  flags, and its return is concentrated at the data-screening level rather than
  the order-override level, which rationalizes the opposing signs of the two
  field experiments on human discretion
  \citep{kesavan2025judgmental,kesavan2020override}.

  \item \textbf{Numerical evidence.} On a calibrated instance with normal demand
  we confirm the unique equilibrium, the exact exposure law, and the simulated
  order path against the mean-field prediction across critical ratios,
  contamination levels, and horizons, and we report the items most exposed to the
  loop (Section~\ref{sec:num}).
\end{enumerate}

% P6 Roadmap.
Section~\ref{sec:lit} places the result in the literature. Section~\ref{sec:model}
states the model and the learner. Section~\ref{sec:law} derives the equilibrium,
the exposure law, and the value of the filter. Section~\ref{sec:num} reports the
numerical study.

\section{Related Literature}\label{sec:lit}

Our work draws on six strands, and we keep to the top journals in each.

\paragraph{Information stalking and the exploration trap.}
A censored learner that only exploits under-orders and so never observes the upper
tail, converging to a level below the optimum
\citep{lariviere1999stalking,ding2002censored}. The known remedy is to order above
the myopic level to acquire information. Our equilibrium looks similar but is
driven by contamination rather than by myopia, and we show that the remedy of
de-censoring does not remove it, because masked stockouts suppress the labels the
de-censoring relies on.

\paragraph{Adaptive learning and regret under censoring.}
A mature line designs censoring-aware policies with vanishing regret, through
distribution-free stochastic-gradient learning
\citep{godfrey2001adaptive,huh2009nonparametric}, the Kaplan--Meier estimator
\citep{huh2011adaptive}, the isolation of the regret cost of censoring
\citep{besbes2013implications}, inference for the resulting policies
\citep{ban2020confidence}, and the recent offline analysis of
\citet{hssaine2024datadriven} that makes the observable boundary explicit. We
adopt this regret language and study the realistic learner that is censoring-aware
yet contamination-blind.

\paragraph{Spiral-down and decision-dependent demand.}
\citet{cooper2006spiral} formalize the loop in which censored observations feed
forecasts that lower controls that further censor observations. Our descent is a
structured instance of performative prediction, namely the setting where deploying
a model shifts the distribution it then faces \citep{perdomo2020performative}. We
give the operations form, derive its equilibrium in closed form, and tie it to
oversight design.

\paragraph{Data-driven, machine learning, and reinforcement learning newsvendor.}
The feature-based data-driven newsvendor \citep{ban2019bigdata}, prescriptive
analytics \citep{bertsimas2020predictive,elmachtoub2022smart}, censoring-aware
feature-based control \citep{ding2024feature}, learning under shifting demand
\citep{chen2021shifting}, and deep reinforcement learning for lost-sales inventory
\citep{gijsbrechts2022deeprl} all engineer performance under censoring. Our point
is orthogonal. Even a well-engineered censoring-aware learner is driven below
target by masked stockouts absent a human filter.

\paragraph{Inventory record inaccuracy and masked stockouts.}
\citet{dehoratius2008empirical,dehoratius2008retail} document that retail systems
often carry phantom inventory, that is, positive on-hand records when the shelf is
empty. \citet{mersereau2015censored,chen2015visibility} show that this inaccuracy
compounds the downward bias already present under demand censoring. We route the
same force through a repeated learning loop in which the identification boundary is
the order itself, and tie the remedy to the returns-and-corrections log that records
when offsets are booked \citep{kesavan2025judgmental}.

\paragraph{Robust and distribution-free inventory.}
The contamination model is in the spirit of Huber's
\citeyearpar{huber1964robust} $\epsilon$-contamination, the classical robustness
device, and connects to distribution-free inventory in the tradition of
\citet{scarf1958minmax}. We are not aware of prior work routing contamination
through a censored, repeated inventory loop in which the identification boundary is
itself a decision.

\paragraph{Human-in-the-loop and algorithm aversion.}
Field experiments show that human override helps at the data-input level and hurts
at the decision level \citep{kesavan2025judgmental,kesavan2020override}, a
behavioral literature documents algorithm aversion and how modest control over the
algorithm restores adoption \citep{dietvorst2018overcoming}, and related work asks
how to elicit and combine human judgment \citep{ibrahim2021eliciting}. We give a
model in which the human's value is precisely as a data filter, which rations the
scarce attention these papers measure.

\section{Model and the Autonomous Learner}\label{sec:model}

Demand $D_t$ is drawn i.i.d.\ from a continuous, strictly increasing distribution
$F$ with density $f$. The newsvendor underage and overage costs are $b$ and $h$,
the critical ratio is $\beta=b/(b+h)$, and the cost-minimizing order is the
$\beta$-quantile $q^\star=F^{-1}(\beta)$. Each period the firm orders $q_t$ and
observes sales and a stockout label.

\paragraph{Contamination through returns and corrections.}
In each period, returns and data corrections are logged as an operational byproduct,
independent of demand. With intensity $\epsilon$, a true high-demand event (demand
$D_t \ge q_t$ that should be recorded as a stockout) coincides with a return or
correction being booked to that period, reducing the recorded demand and masking
the stockout. Specifically, the recorded demand becomes a value drawn from a
downward source $G$, where $G$ is stochastically dominated by $F$ and represents
the magnitude of typical returns or corrections. The recorded value and the
operative stockout status are now decoupled, and the stockout label follows the
recorded value rather than the true demand. This is the one error a machine cannot
correct on its own without domain knowledge or access to the returns and corrections
log, which contains the ground truth for when this offset occurred.

\paragraph{Recorded data.}
Work on a probability space $(\Omega,\mathcal F,\mathbb P)$ with a filtration
$(\mathcal F_t)$, and let the order $q_t$ be $\mathcal F_{t-1}$-measurable. In period
$t$, given $q_t$, nature draws the demand $D_t\sim F$, a contamination coin
$U_t\sim\mathrm{Unif}[0,1]$, and if contamination occurs, a correction value $D_t^G\sim G$,
mutually independent and independent of $\mathcal F_{t-1}$. The recorded outcome is one of
three. If $D_t<q_t$ the day is served and recorded at value $D_t$. If $D_t\ge q_t$
and $U_t>\epsilon$ the day is a genuine high-demand event, recorded as censored at $q_t$
(stockout). If $D_t\ge q_t$ and $U_t\le\epsilon$ a correction is booked for that period,
reducing recorded demand to $D_t^G$, which enters the learner's data as an uncensored
low-sales observation. The returns-and-corrections log records which periods experienced
a booking; call this a \textit{logged correction day}. Let
$g_t=\mathbf 1\{\text{recorded demand}<q_t\}$ denote the recorded served-indicator,
which is $\mathcal F_t$-measurable.

\paragraph{The learner and the audit.}
The autonomous policy is censoring-aware and contamination-blind. It uses the
stockout labels, so it scores a genuine high-demand event as demand above the order,
but it cannot tell a masked high-demand day from a true low-sales day without
access to the returns-and-corrections log. We analyze the canonical adaptive form,
a projected Robbins--Monro recursion that drives the recorded served-rate to the
target $\beta$,
\begin{equation}\label{eq:rm}
  q_{t+1}=\Pi_{[\underline q,\overline q]}\!\big(q_t-a_t\,(g_t-\beta)\big),
  \qquad a_t>0,\quad \textstyle\sum_t a_t=\infty,\quad \sum_t a_t^2<\infty,
\end{equation}
where $\Pi$ is the projection onto a compact interval $[\underline q,\overline q]$
that contains $q^\star$ in its interior. The estimate-then-optimize learner that
de-censors by the Kaplan--Meier estimator over a rolling window and orders the
$\beta$-quantile of its estimate is an alternative with the same mean field, by
Lemma~\ref{lem:mf}, and is the version we simulate in Section~\ref{sec:num}. A human
filter operates at the data level by consulting the returns-and-corrections log and,
for a fraction $\varphi$ of the logged corrections, undoing the recorded offset to
recover the true demand signal and feeding it back into the learner's data stream.
This data-level human intervention is shown to improve profit in the field, in
contrast to order-level overrides which reduce it \citep{kesavan2025judgmental}.

\section{The Contamination Spiral and the Exposure Law}\label{sec:law}

We now ask where the loop settles, how fast it descends, and what restores it,
keeping the clean benchmark $q^\star$ in view throughout. The first step records
what the learner perceives at a fixed order and reduces the recursion to a
one-dimensional fixed-point problem.

\begin{lemma}[Perceived demand and the mean field]\label{lem:mf}
Fix an order $q\in[\underline q,\overline q]$.
\textup{(a)} Operating persistently at $q$, the estimate-then-optimize learner's
windowed estimate converges almost surely and uniformly to
\begin{equation}\label{eq:Fhat}
  \widehat F_q(x)=F(x)+\epsilon\,(1-F(q))\,G(x),\qquad x\le q .
\end{equation}
\textup{(b)} The recorded served-indicator satisfies
$\mathbb E[g_t\mid q_t=q]=\widehat F_q(q)=H(q)$, where
\begin{equation}\label{eq:H}
  H(q)\;:=\;F(q)+\epsilon\,(1-F(q))\,G(q).
\end{equation}
\end{lemma}

\noindent\textit{Proof.}
(a) At the fixed level $q$ the recorded values are i.i.d., and with a single
censoring level the Kaplan--Meier estimator equals the empirical distribution on
$[0,q]$ with the censored mass placed above $q$. The uncensored records at or below
$x<q$ are the genuine demands below $x$, of probability $F(x)$, together with the
masked stockouts whose correction value falls below $x$, of probability
$\epsilon(1-F(q))G(x)$. Glivenko--Cantelli gives almost sure uniform convergence to
their sum \eqref{eq:Fhat}. (b) The indicator $g_t$ equals one on exactly the served
and masked-below days, so
$\mathbb E[g_t\mid q_t=q]=F(q)+\epsilon(1-F(q))G(q)=\widehat F_q(q)$. \hfill$\square$

The genuine uncensored days below $x$ contribute $F(x)$. The masked stockout
days contribute the second term, in which $\epsilon(1-F(q))$ is the share of days
that stock out and are masked in the feed. The factor $1-F(q)$ is the stockout
frequency,
and it ties the size of the contamination to the order itself. That single coupling
is what turns a static bias into a loop, because the equilibrium order solves
\begin{equation}\label{eq:fp}
  H(q)=\beta .
\end{equation}
The whole paper rests on the following result, which reads the loop in three parts,
namely where it goes, how much it costs at the margin, and how cheaply a human
reverses it.

\begin{mdframed}[linewidth=0.6pt,roundcorner=3pt,backgroundcolor=black!2]
\begin{theorem}[Contamination spiral and the exposure law]\label{thm:main}
Let $F$ be continuous and strictly increasing with $q^\star=F^{-1}(\beta)$,
$\beta\in(0,1)$, and let $G$ be stochastically dominated by $F$. Under the
censoring-aware, contamination-blind learner of Section~\ref{sec:model} with
contamination intensity $\epsilon$, the following hold.
\begin{enumerate}[label=\textup{(\roman*)},leftmargin=2em,itemsep=0.3em]
  \item \textup{(Spiral to a unique equilibrium, no tipping point)} The map $H$ is
  strictly increasing, so \eqref{eq:fp} has a unique root $q_\infty(\epsilon,\beta)$,
  and the projected recursion \eqref{eq:rm} converges to it almost surely from any
  start in $[\underline q,\overline q]$. The root satisfies $q_\infty=q^\star$ at
  $\epsilon=0$ and $q_\infty<q^\star$, strictly decreasing in $\epsilon$, for
  $\epsilon>0$. No second equilibrium or tipping point exists.

  \item \textup{(Exposure law)} The equilibrium service level $F(q_\infty)$ is
  smooth and strictly decreasing in $\epsilon$, with
  \begin{equation}\label{eq:law}
    \frac{\mathrm{d}\,F(q_\infty)}{\mathrm{d}\epsilon}\bigg|_{\epsilon=0}
    =-(1-\beta)\,G(q^\star),
  \end{equation}
  which equals $-(1-\beta)$ when the contamination lies below $q^\star$. The
  marginal service loss is the item's stockout frequency.

  \item \textup{(Cost and oversight)} Long-run average cost exceeds the optimum by
  \begin{equation}\label{eq:cost}
    C(q_\infty)-C(q^\star)
    =\tfrac12(b+h)f(q^\star)\,(q^\star-q_\infty)^2+o\big((q^\star-q_\infty)^2\big)
    =\Theta\!\Big(\tfrac{(b+h)\,\epsilon^2(1-\beta)^2}{f(q^\star)}\Big).
  \end{equation}
  A correction-log alarm that audits a fraction $\varphi$ of logged correction
  days corrects the masked records at that rate, reducing the effective
  contamination to $\epsilon(1-\varphi)$ and the cost penalty to a
  $(1-\varphi)^2$ fraction, so half coverage removes three quarters of the penalty
  and full coverage restores $q^\star$. The value of oversight is largest for
  low-service items.
\end{enumerate}
\end{theorem}
\end{mdframed}

\noindent\textit{Proof of Theorem~\ref{thm:main}.}
Part (i), structure. Differentiating \eqref{eq:H},
$H'(q)=f(q)\,(1-\epsilon G(q))+\epsilon\,(1-F(q))\,g(q)$, and both terms are
nonnegative because $\epsilon G(q)\le1$, with the first strictly positive, so $H$
is strictly increasing and the root of \eqref{eq:fp} is unique. The added term in
\eqref{eq:H} is positive for $\epsilon>0$, so $F(q_\infty)<\beta$, that is
$q_\infty<q^\star$, and monotonicity in $\epsilon$ follows from the implicit
function theorem.

Part (i), convergence. Write $h(q)=\beta-H(q)$. By Lemma~\ref{lem:mf}(b) the
recursion \eqref{eq:rm} is $q_{t+1}=\Pi_{[\underline q,\overline q]}(q_t+a_t h(q_t)
+a_t M_t)$ with $M_t=\mathbb E[g_t\mid\mathcal F_{t-1}]-g_t$, a martingale-difference
sequence with $|M_t|\le1$. Strict monotonicity of $H$ gives the Lyapunov sign
condition $(q-q_\infty)h(q)<0$ for every $q\ne q_\infty$ in
$[\underline q,\overline q]$. Set $V_t=(q_t-q_\infty)^2$. Since $q_\infty$ lies in
the interior and $\Pi$ is nonexpansive with $\Pi(q_\infty)=q_\infty$,
\[
  \mathbb E[V_{t+1}\mid\mathcal F_{t-1}]
  \;\le\; V_t+2a_t\,(q_t-q_\infty)\,h(q_t)+a_t^2\,\mathbb E[M_t^2\mid\mathcal F_{t-1}]
  \;\le\; V_t-2a_t\,\Delta_t+a_t^2,
\]
where $\Delta_t=(q_t-q_\infty)(H(q_t)-\beta)\ge0$. As $\sum_t a_t^2<\infty$, the
Robbins--Siegmund nonnegative almost-supermartingale theorem
\citep{robbins1971convergence} implies that $V_t$ converges almost surely and
$\sum_t a_t\Delta_t<\infty$ almost surely. Were $|q_t-q_\infty|\ge\delta$ for all
large $t$ on a positive-probability set, then $\Delta_t$ would be bounded below by
$\inf_{\delta\le|q-q_\infty|}\Delta(q)>0$ there, forcing $\sum_t a_t\Delta_t=\infty$
because $\sum_t a_t=\infty$, a contradiction. Hence $\liminf_t|q_t-q_\infty|=0$, and
since $V_t$ converges, $q_t\to q_\infty$ almost surely
\citep{robbins1951stochastic,kushner2003stochastic}. This holds from any start in
$[\underline q,\overline q]$, so there is no second equilibrium and no tipping
point.

Part (ii). Differentiate $H(q_\infty(\epsilon))=\beta$ in $\epsilon$ at
$\epsilon=0$, where $q_\infty=q^\star$ and $F(q^\star)=\beta$. The cross terms in
$\epsilon$ vanish, leaving $f(q^\star)\,q_\infty'(0)+(1-\beta)\,G(q^\star)=0$, and
multiplying by $f(q^\star)$ gives \eqref{eq:law}.
Part (iii). Since $C'(q)=(b+h)F(q)-b$, we have $C'(q^\star)=0$ and
$C''(q^\star)=(b+h)f(q^\star)$, which gives the quadratic expansion. Substituting
$q^\star-q_\infty=\epsilon(1-\beta)G(q^\star)/f(q^\star)+o(\epsilon)$ from
\eqref{eq:law} gives the order term. A correction-log alarm reviews a fraction
$\varphi$ of logged correction days. Each masked stockout appears in that log, so
an audited day is corrected with probability $\varphi$, the effective intensity is
$\epsilon(1-\varphi)$, and the cost, quadratic in intensity by \eqref{eq:cost},
scales as $(1-\varphi)^2$. \hfill$\square$

The same structure pins down the rate, and with it the cumulative cost behind the
$\Theta(T)$ claim. The step size enters through the local slope
$\lambda:=H'(q_\infty)>0$ established in part (i).

\begin{lemma}[Rate]\label{lem:rate}
Under \eqref{eq:rm} with $a_t=a/t$ and $2a\lambda>1$,
$\mathbb E[(q_t-q_\infty)^2]=O(1/t)$ and $\bigl|\mathbb E[q_t-q_\infty]\bigr|=O(1/t)$.
\end{lemma}

\noindent This is the classical stochastic-approximation rate for a recursion whose
mean field is locally linear and strictly stable at its root
\citep{kushner2003stochastic}. It controls the transient and isolates the
contamination bias $q^\star-q_\infty$ as the term that persists.

\begin{lemma}[Cumulative regret]\label{lem:regret}
Let $R_T=\sum_{t=1}^{T}\bigl(\mathbb E[C(q_t)]-C(q^\star)\bigr)$ under the step size
of Lemma~\ref{lem:rate}. Then
\[
  R_T=T\,\bigl(C(q_\infty)-C(q^\star)\bigr)+O(\log T)
     =\Theta\!\Big(\tfrac{(b+h)\,\epsilon^2(1-\beta)^2}{f(q^\star)}\;T\Big).
\]
\end{lemma}

\noindent\textit{Proof.}
A second-order expansion gives
$C(q_t)-C(q_\infty)=C'(q_\infty)(q_t-q_\infty)+\tfrac12C''(\zeta_t)(q_t-q_\infty)^2$
for some $\zeta_t$ between $q_t$ and $q_\infty$, with $C''$ bounded on
$[\underline q,\overline q]$. Taking expectations and applying Lemma~\ref{lem:rate}
to both terms gives $\mathbb E[C(q_t)]-C(q_\infty)=O(1/t)$, which sums to
$O(\log T)$. Adding $T\,(C(q_\infty)-C(q^\star))$ and substituting the constant
order from \eqref{eq:cost} gives the claim. \hfill$\square$

Theorem~\ref{thm:main} captures an item run unattended on censored, contaminated
data, and the mechanism is the factor $1-F(q)$ in \eqref{eq:Fhat}. The learner
cannot separate a genuinely weak demand signal from a corrupted one, its
censoring-aware response to an apparently weak signal is to order less, and
ordering less raises the stockout frequency that feeds the next period more
corrupted records. Two readings follow. First, the descent has no warning sign,
because $H$ is monotone and there is no second equilibrium, so a manager watching
for a sudden break will see none and the shortfall arrives instead as a slow,
permanent erosion. Second, the marginal damage is the stockout frequency $1-\beta$,
so a firm can rank its catalogue by target service level with data already in hand
and find that the cheap, fast-moving, thin-service items it automates first are
exactly the items the loop punishes hardest. The failure mode to remember is that
de-censoring does not protect the item and a better optimizer does not help,
because the data are the problem. The complementary priority is a human who screens
logged correction days where the masking enters, and because the penalty is
quadratic in contamination, auditing even half of those days removes most of it.

\section{Numerical Study}\label{sec:num}

We calibrate demand as $F=\mathcal N(100,30)$ and the contamination source as
$G=\mathcal N(40,10)$, that is returns batches booked at about $40$ units. We
simulate the learner with $N=200$ demand draws per period, a rolling window of
$W=60$ periods, horizon $T=4000$, and a generous start of $q_0=160$, and we add a
correction-log alarm that audits a fraction $\varphi$ of logged correction days.
Table~\ref{tab:sim} reports the long-run order against the mean-field equilibrium,
the realized service, the cost gap above the optimum, and the audit load, namely
the fraction of item-days the human inspects.

\begin{table}[t]
\centering
\small
\caption{Simulated learner and correction-log filter. Demand $\mathcal N(100,30)$,
contamination $G=\mathcal N(40,10)$ on masked stockout days, $\epsilon=0.30$ for the
contaminated rows. $\varphi$ is the fraction of logged correction days audited. The
optimum is $q^\star=84.3$ for $\beta=0.30$ and $q^\star=115.7$ for $\beta=0.70$.
Cost gap is percent above optimal newsvendor cost. Audit load is the fraction of
item-days inspected. Standard deviations are over the final $500$ periods.}
\label{tab:sim}
\begin{tabular}{c c c c c c c c}
\toprule
$\beta$ & $\epsilon$ & $\varphi$ & mean-field $q_\infty$ & sim $q_\infty$ (s.d.)
 & service & cost gap \% & audit load\\
\midrule
0.30 & 0.00 & --  & 84.3 & 83.2 \,(0.02) & 0.287 & 0.1  & 0.000\\
0.30 & 0.30 & 0.0 & 51.3 & 51.6 \,(0.61) & 0.053 & 45.7 & 0.000\\
0.30 & 0.30 & 0.5 & 72.1 & 72.5 \,(0.65) & 0.180 & 7.1  & 0.122\\
0.30 & 0.30 & 1.0 & 84.3 & 83.1 \,(0.04) & 0.287 & 0.1  & 0.212\\
\midrule
0.70 & 0.00 & --  & 115.7 & 114.2 (0.02) & 0.682 & 0.1  & 0.000\\
0.70 & 0.30 & 0.0 & 105.4 & 104.2 (0.03) & 0.555 & 7.8  & 0.000\\
0.70 & 0.30 & 0.5 & 111.3 & 110.1 (0.03) & 0.631 & 1.8  & 0.054\\
0.70 & 0.30 & 1.0 & 115.7 & 114.3 (0.04) & 0.683 & 0.1  & 0.093\\
\bottomrule
\end{tabular}
\end{table}

We simulate the estimate-then-optimize learner, which shares the mean field of the
analyzed Robbins--Monro recursion by Lemma~\ref{lem:mf}. The Robbins--Monro learner
of \eqref{eq:rm} converges to the same equilibria, matching the roots of
\eqref{eq:fp} to within $0.2$ units across the cases below.

The table confirms every part of Theorem~\ref{thm:main}. The unfiltered learner
($\varphi=0$) starts well above the optimum and descends to the mean-field
equilibrium to within about a unit, validating part~(i) and the absence of a
tipping point. The damage is large and ranked by service level as part~(ii)
predicts. The low-service item at $\beta=0.30$ loses almost all of its service and
runs $46$ percent above optimal cost, while the high-service item at $\beta=0.70$
loses far less. The filter validates part~(iii) and its convex recovery. Auditing
half of the logged correction days removes more than three quarters of the cost
penalty, from $45.7$ to $7.1$ percent at $\beta=0.30$ and from $7.8$ to $1.8$
percent at $\beta=0.70$, and full coverage restores the optimum. The audit load
shows the trade-off the manager faces. Because the alarm fires only on logged
correction days, namely the masked stockouts themselves, full recovery requires
reviewing about $\epsilon(1-F(q_\infty))$ of item-days, far less than the stockout
rate the old stockout-triggered alarm would have implied, and half coverage already
captures most of the value.

\bibliography{references}

\end{document}
